% !TeX root = ../tesis.tex
%%%%%%%%%%%%%%%%%%%%%%%%%%%%%%%%%%%%%%%%%%%%%%%%%%%%%%%%%%%%%%%%%%%
% CAPÍTULO 5 — ANÁLISIS DE LA RED DE TRANSPORTE ACTUAL DE LA ZMVM
%
% Issue #24 [C-7]: Reestructuración del capítulo de resultados
% como diagnóstico territorial de la red actual.
%
% PREREQUISITO: Verificar notebooks VFTModel antes de redactar.
% Ver: /VFTModel/notebooks/ (fuera del repo)
%
% Estructura objetivo:
%   §5.1  La Red Metropolitana como Objeto de Análisis
%   §5.2  Construcción del Grafo Topológico
%   §5.3  Indicadores Topológicos — Fase 1
%         §5.3.1  Validación de Impedancia
%         §5.3.2  Nivel de Cobertura Espacial (C)
%         §5.3.3  Fuerza Capilar (FC_h)
%         §5.3.4  Factor de Desviación (DI)
%   §5.4  Ponderación Dinámica — Fase 2
%         §5.4.1  Coeficiente de Fricción Vial (CF)
%         §5.4.2  Penalización por Transferencia (Tb)
%   §5.5  Análisis de Red — Fase 3
%         §5.5.1  Tiempo de Viaje Promedio (T)
%         §5.5.2  Centralidad de Intermediación B(v)
%         §5.5.3  Vulnerabilidad Estructural ΔE
%   §5.6  Diagnóstico Territorial — Zonas de Déficit (NUEVO)
%         §5.6.1  Cobertura por cuadrante del Anillo Periférico
%         §5.6.2  Zonas con mayor déficit estructural
%         §5.6.3  Nodos críticos → enlace a Cap. 6
%%%%%%%%%%%%%%%%%%%%%%%%%%%%%%%%%%%%%%%%%%%%%%%%%%%%%%%%%%%%%%%%%%%
\chapter{Análisis de la Red de Transporte Actual de la \zmvm}
\label{ch:analisis-red-actual}

% Secciones del capítulo (se agregan conforme se redacten)
% !TeX root = ../tesis.tex
%%%%%%%%%%%%%%%%%%%%%%%%%%%%%%%%%%%%%%%%%%%%%%%%%%%%%%%%%%%%%%%%%%%
% SECCIÓN 5.1 — LA RED METROPOLITANA COMO OBJETO DE ANÁLISIS
% Issue #24 [C-7]
%%%%%%%%%%%%%%%%%%%%%%%%%%%%%%%%%%%%%%%%%%%%%%%%%%%%%%%%%%%%%%%%%%%
\section{La Red Metropolitana como Objeto de Análisis}
\label{sec:red-objeto-analisis}

El diagnóstico del sistema de transporte público de la \zmvm{} enfrenta un
problema de escala: la red integra once modalidades de transporte masivo
distribuidas sobre un territorio metropolitano que trasciende la \cdmx{} e
incorpora municipios del Estado de México en una continuidad urbana funcional
\citep{semovi2020diagnostico}. Los indicadores de movilidad son difíciles de
aplicar de forma homogénea a cada sistema de transporte, y este problema no
radica en la falta de datos: el Gobierno de la \cdmx{} libera sistemáticamente
la información de sus redes a través del Portal de Datos Abiertos
\citep{adip_gtfs_2024}. La barrera real es la heterogeneidad estructural de
esos datos: once sistemas con esquemas, frecuencias, geometrías y atributos
distintos que, sin una herramienta capaz de integrarlos bajo un modelo común,
no pueden analizarse como red sino solo como inventario.

Este es el problema que el \textbf{Apímetro} y el \textbf{Modelo VFT},
desarrollados y documentados en el Capítulo~4, están diseñados para resolver
\citep{cabrera2025apimetro, cabrera2026VFTModel}. Al representar la red como
un grafo dirigido y ponderado —donde cada estación es un nodo y cada segmento
operativo una arista con peso temporal— los once sistemas quedan integrados
bajo un modelo matemático común \citep{newman2010networks}. Sobre esa
representación es posible calcular propiedades estructurales independientes de
la demanda: la cobertura territorial de la red, la concentración de flujos en
los nodos de transferencia, la eficiencia de los trayectos y la fragilidad de
los puntos críticos.

El presente capítulo aplica esa herramienta al sistema de transporte actual de
la \zmvm{} para responder tres preguntas estructurales: ¿qué territorio cubre
la red y dónde se concentran los flujos? ¿Cuánto tarda en promedio
desplazarse entre cualquier par de estaciones y qué nodos articulan la mayor
parte de esos trayectos? ¿Qué tan frágil es la red ante la pérdida de sus
nodos más críticos? Las respuestas se organizan en tres fases de análisis
—indicadores topológicos, ponderación dinámica y análisis de red— y culminan
en un diagnóstico territorial que identifica las zonas de mayor déficit
estructural como punto de partida para el corredor anillar hipotético del
Capítulo~6.

El Cuadro~\ref{tab:resumen-indicadores-cap5} sintetiza los ocho indicadores
que componen el análisis, su fase de cálculo y su estado de avance.

\begin{table}[H]
    \centering
    \caption[Indicadores del Modelo VFT por fase de cálculo]{Resumen de los
    indicadores del \textbf{Modelo VFT}, su fase de cálculo, dependencias y
    estado de avance en esta investigación.}
    \label{tab:resumen-indicadores-cap5}
    \begin{tabularx}{\textwidth}{|l|c|X|l|}
        \hline
        \textbf{Indicador} & \textbf{Fase} & \textbf{Dependencias} &
        \textbf{Estado} \\
        \hline
        Validación de Impedancia      & 1 & Grafo base + fallbacks       & Calculado \\
        Cobertura Espacial ($C$)      & 1 & Geometría de estaciones      & Calculado \\
        Fuerza Capilar ($FC$)         & 1 & Grafo base                   & Calculado \\
        Factor de Desviación ($DI$)   & 1 & Grafo base + impedancia      & Calculado \\
        \hline
        Coef. de Fricción Vial ($CF$)     & 2 & Fase 1 completa          & Calculado \\
        Penalización por Transf. ($T_b$)  & 2 & Fase 1 + $CF$            & Calculado \\
        Node Score ($\text{VFT\_score}$)  & 2 & Fase 1 + datos MV1/MV2   & En diseño \\
        \hline
        Tiempo de Viaje Prom. ($T$)   & 3 & Fases 1 y 2 completas        & Calculado \\
        Centralidad Interm. ($B$)     & 3 & Fases 1 y 2 + $T$            & Calculado \\
        Vulnerabilidad ($\Delta E$)   & 3 & Fases 1, 2 y $B$             & Calculado \\
        \hline
    \end{tabularx}
    \fuentefigura{Elaboración propia con base en \texttt{VFTModel}
    \citep{cabrera2026VFTModel}.}
\end{table}

% !TeX root = ../tesis.tex
%%%%%%%%%%%%%%%%%%%%%%%%%%%%%%%%%%%%%%%%%%%%%%%%%%%%%%%%%%%%%%%%%%%
% SECCIÓN 5.2 — CONSTRUCCIÓN DEL GRAFO TOPOLÓGICO
% Issue #24 [C-7]
%%%%%%%%%%%%%%%%%%%%%%%%%%%%%%%%%%%%%%%%%%%%%%%%%%%%%%%%%%%%%%%%%%%
\section{Construcción del Grafo Topológico}
\label{sec:grafo-topologico}

A partir de los datos geoespaciales procesados por el \textbf{Apímetro}, el
\textbf{Modelo VFT} construyó un grafo dirigido y ponderado que representa la
red de transporte masivo de la \zmvm{} como objeto de análisis estructural
(véanse Capítulo~4 y Anexo~B para el detalle de implementación). El proceso de
integración topológica —mediante un árbol \texttt{KDTree} con tolerancia de
85~m, corrida del 18 de abril de 2026— redujo las 22{,}766 entidades espaciales
del dataset de origen a \textbf{11{,}115 nodos canónicos} conectados por
\textbf{45{,}679 aristas}, de las cuales 25{,}197 corresponden a segmentos de
servicio real y 20{,}482 a aristas de transbordo peatonal generadas
algorítmicamente \citep{cabrera2026VFTModel}.

\begin{figure}[H]
    \centering
    % Diagrama: Pipeline de procesamiento interno del Modelo VFT (4 etapas)
% Usado en: 5-Resultados/5-1-Introduccion-Resultados.tex  (fig:arq-capas)
% Reemplaza: 3 subfiguras arq_capa1/2/3.png con captions incorrectos
\begin{tikzpicture}[
    node distance = 0.6cm and 0.5cm,
    etapa/.style = {
        rectangle, draw=unamAzul, fill=unamAzul!8,
        text width=2.8cm, align=center,
        minimum height=1.5cm, rounded corners=4pt,
        font=\small\bfseries
    },
    dato/.style = {
        font=\scriptsize\itshape, color=unamAzul!75,
        align=center, text width=2.8cm
    },
    cifra/.style = {
        font=\footnotesize\bfseries, color=unamAzul,
        align=center, text width=2.8cm
    },
    flecha/.style = {-Stealth, thick, color=unamAzul!65}
]

%% — Etapas —
\node[etapa] (ing) {Ingesta\\GTFS};
\node[etapa, right=of ing]  (snap) {Snapping\\Lógico};
\node[etapa, right=of snap] (grf)  {Construcción\\del Grafo};
\node[etapa, right=of grf]  (imp)  {Ponderación\\Impedancia};

%% — Cifras sobre cada etapa —
\node[cifra, above=0.18cm of ing]  {22{,}766 entidades};
\node[cifra, above=0.18cm of snap] {11{,}115 nodos};
\node[cifra, above=0.18cm of grf]  {45{,}679 aristas};
\node[cifra, above=0.18cm of imp]  {grafo $G$ dirigido};

%% — Detalles bajo cada etapa —
\node[dato, below=0.18cm of ing]  {Apimetro REST\\GeoJSON};
\node[dato, below=0.18cm of snap] {KDTree · $d \leq 50$~m\\eliminación duplicados};
\node[dato, below=0.18cm of grf]  {NetworkX\\25{,}197 servicio\\20{,}482 transbordo};
\node[dato, below=0.18cm of imp]  {Motor VFT\\pesos en minutos};

%% — Flechas —
\draw[flecha] (ing)  -- (snap);
\draw[flecha] (snap) -- (grf);
\draw[flecha] (grf)  -- (imp);

%% — Nota conceptual: abstracción Red Física → Grafo Matemático —
\node[font=\scriptsize\color{gray!70}, anchor=north]
    at ($(ing.south)!0.5!(imp.south) + (0,-1.15cm)$)
    {Red Física $\longrightarrow$ Grafo Matemático $\mathcal{G}(\mathcal{V},\,\mathcal{E},\,w)$};

\end{tikzpicture}

    \caption[Pipeline de procesamiento del Modelo VFT]{Las cuatro etapas que
    transforman los datos brutos de \textbf{Apímetro} en el grafo topológico
    operativo. El \termino{Snapping Lógico Geográfico} mediante árbol
    \texttt{KDTree} (tolerancia: 85~m) reduce las 22{,}766 entidades originales
    a 11{,}115 nodos canónicos; la ponderación final asigna impedancias
    temporales a las 45{,}679 aristas del grafo dirigido
    $\mathcal{G}(\mathcal{V},\mathcal{E},w)$.}
    \label{fig:arq-capas}
    \fuentefigura{Elaboración propia con \texttt{VFTModel}
    \citep{cabrera2026VFTModel}.}
\end{figure}

\begin{itemize}
    \item \textbf{11{,}115 nodos topológicos} — estaciones operativas tras
    eliminar nodos fantasma y redundancias geométricas.
    \item \textbf{45{,}679 aristas}: 25{,}197 segmentos de servicio real y
    20{,}482 aristas de transbordo peatonal generadas por el
    \termino{Snapping Lógico Geográfico}.
    \item \textbf{11 sistemas de transporte} representados: \textsc{Metro},
    Metrobús, Trolebús, Tren Ligero, Tren Suburbano, Cablebús, Mexicable,
    Mexibús, \textsc{rtp}, Corredores Concesionados y Pumabús.
    \item Distribución modal: \textbf{12\,\% infraestructura estructurada} con
    derecho de vía exclusivo o confinado (\textsc{Metro}, Metrobús, Tren Ligero,
    Cablebús, Tren Suburbano) y \textbf{88\,\% transporte de superficie} en
    vía compartida o mixta (\textsc{rtp}, Corredores Concesionados, Trolebús,
    Mexibús, Mexicable, Pumabús).
    \item \textbf{Componente gigante}: 10{,}561 nodos (95\,\% del total),
    sobre la cual se ejecutan los algoritmos de Fase~3.
\end{itemize}

El Cuadro~\ref{tab:aristas-por-sistema} detalla la distribución de aristas de
servicio real por sistema, el tipo de derecho de vía asignado y el coeficiente
de fricción vial correspondiente.

\begin{table}[H]
    \centering
    \caption[Distribución de aristas de servicio por sistema de transporte]{
    Distribución de aristas de servicio real por sistema de transporte en el
    grafo de la \zmvm{}. El coeficiente $CF$ refleja el tipo de derecho de vía.
    Corrida: 2026-04-18.}
    \label{tab:aristas-por-sistema}
    \begin{tabular}{|l|r|l|c|}
        \hline
        \textbf{Sistema} & \textbf{Aristas} & \textbf{Derecho de vía} &
        \textbf{$CF$} \\
        \hline
        \textsc{rtp}                  & 11{,}041 & Mixto       & 1.759 \\
        Corredores Concesionados (CC) & 10{,}603 & Mixto       & 1.759 \\
        Trolebús\textsuperscript{†}   &  1{,}506 & Compartido  & 1.380 \\
        Metrobús (MB)                 &    792   & Exclusivo   & 1.000 \\
        \textsc{Metro} (STC)          &    528   & Exclusivo   & 1.000 \\
        Mexibús\textsuperscript{†}    &    391   & Confinado   & 1.152 \\
        Pumabús                       &    234   & Mixto       & 1.759 \\
        Tren Ligero (TL)              &     52   & Exclusivo   & 1.000 \\
        Cablebús (CBB)\textsuperscript{‡} &  30  & Confinado   & 1.152 \\
        Mexicable\textsuperscript{†}  &     17   & Confinado   & 1.152 \\
        Tren Suburbano (SUB)          &      3   & Exclusivo   & 1.000 \\
        \hline
        \textbf{Total}                & \textbf{25{,}197} & & \\
        \hline
    \end{tabular}
    \vspace{0.5em}

    \noindent\textsuperscript{†}\,A la fecha de ingesta (2026-04-18), la
    Línea 13 del Trolebús, el Mexicable y el Mexibús no contaban con registros
    publicados en el Portal de Datos Abiertos de la \cdmx{}. Su geometría y
    paradas fueron integradas a partir de datos de
    \termino{OpenStreetMap} \citep{osm2012}.

    \noindent\textsuperscript{‡}\,El sistema Cablebús registra únicamente
    30 aristas, cifra inferior a la esperada para sus tres líneas operativas.
    Esta subrepresentación es una limitación del \textit{backend} del
    \textbf{Apímetro} en la fecha de corrida y no afecta los indicadores de
    los sistemas restantes.

    \fuentefigura{Elaboración propia con \texttt{VFTModel}
    \citep{cabrera2026VFTModel}; datos: \citep{adip_gtfs_2024, osm2012}.}
\end{table}

\begin{figure}[H]
    \centering
    \includegraphics[width=0.92\textwidth]{%
        Figures/Mapas/Qgis_Resources/RedActual/Mapa_Transporte_Publico}
    \caption[Red multimodal de la \zmvm{} por sistema de transporte]{Red
    multimodal de la \zmvm{} organizada por tipo de sistema de transporte,
    elaborada con \textbf{QGIS} a partir de los datos procesados por el
    \textbf{Apímetro}. La versión cartográfica extendida de la red esquemática
    se presenta en la Figura~\ref{fig:anx-mapa-red-esquematica} del
    Anexo~\ref{anx:cartografia-presentacion}.}
    \label{fig:mapa-transporte-publico}
    \fuentefigura{Elaboración propia con \textbf{QGIS};
    datos: \citep{adip_gtfs_2024, osm2012}.}
\end{figure}

El grafo resultante constituye la representación operativa sobre la que se
calculan todos los indicadores presentados en las secciones siguientes. Su
escala —11{,}115 nodos y 45{,}679 aristas distribuidos en once sistemas— hace
de la \zmvm{} uno de los grafos de transporte público más densos documentados
en la literatura latinoamericana \citep{lotero2014vulnerabilidad}, lo que
confiere mayor relevancia a los resultados de los indicadores de cobertura,
centralidad y vulnerabilidad que se exponen a continuación.

% !TeX root = ../../tesis.tex
%%%%%%%%%%%%%%%%%%%%%%%%%%%%%%%%%%%%%%%%%%%%%%%%%%%%%%%%%%%%%%%%%%%
% SECCIÓN 5.3 — INDICADORES TOPOLÓGICOS: FASE 1
% Issue #24 [C-7]
%%%%%%%%%%%%%%%%%%%%%%%%%%%%%%%%%%%%%%%%%%%%%%%%%%%%%%%%%%%%%%%%%%%
\section{Indicadores Topológicos: Fase~1}
\label{sec:fase1-indicadores}

La Fase~1 agrupa los indicadores que caracterizan la red a partir de su
geometría y conectividad, sin requerir datos de demanda ni ponderaciones
dinámicas. Su punto de partida es la validación del motor de impedancia, que
garantiza la coherencia física de los pesos asignados a cada segmento antes
de calcular los indicadores de cobertura, capilaridad y eficiencia de ruta.

% !TeX root = ../../tesis.tex
%%%%%%%%%%%%%%%%%%%%%%%%%%%%%%%%%%%%%%%%%%%%%%%%%%%%%%%%%%%%%%%%%%%
% SECCIÓN 5.3.1 — VALIDACIÓN DEL MOTOR DE IMPEDANCIA
%%%%%%%%%%%%%%%%%%%%%%%%%%%%%%%%%%%%%%%%%%%%%%%%%%%%%%%%%%%%%%%%%%%
\subsection{Validación del Motor de Impedancia Temporal}
\label{subsec:validacion-impedancia}

Para verificar la consistencia del motor de impedancia del \textbf{Modelo VFT},
se calculó la velocidad efectiva implícita de cada sistema tras aplicar el
Coeficiente de Fricción Vial ($CF$) y se contrastó con los rangos operativos
nominales documentados por los organismos operadores de la \zmvm{}. La
velocidad efectiva se obtiene a partir de la
Ecuación~\eqref{eq:velocidad-efectiva}:

\begin{equation}
    V_{\text{efectiva}} = \frac{V_{\text{fallback}}}{CF}
    \label{eq:velocidad-efectiva}
\end{equation}

donde $V_{\text{fallback}}$ representa la velocidad comercial de flujo libre
del sistema y $CF$ es el coeficiente de fricción según el tipo de derecho
de vía (véase §\ref{subsec:fase2-cf}).

El Cuadro~\ref{tab:validacion-impedancia} sintetiza los resultados de la
corrida ejecutada sobre \textbf{25{,}197 segmentos de servicio real}.

\begin{table}[H]
    \centering
    \caption[Validación de velocidades efectivas por sistema de transporte]{
    Contraste entre la velocidad implícita post-fricción calculada por el
    motor VFT y los rangos operativos nominales de cada sistema.
    Corrida: 2026-04-18.}
    \label{tab:validacion-impedancia}
    \begin{tabular}{|l|r|r|c|c|}
        \hline
        \textbf{Sistema} & \textbf{$V$ modelo (km/h)} &
        \textbf{$V$ efectiva (km/h)} &
        \textbf{Rango esperado} & \textbf{Estado} \\
        \hline
        Tren Suburbano   & 65.0 & 65.0 & {[}55--75{]}  & \checkmark\ OK \\
        Metro (STC)      & 36.0 & 36.0 & {[}30--42{]}  & \checkmark\ OK \\
        Tren Ligero      & 22.0 & 22.0 & {[}18--26{]}  & \checkmark\ OK \\
        Trolebús         & 18.0 & 13.0 & {[}14--22{]}  & $\approx$ Límite \\
        Cablebús         & 20.0 & 17.4 & {[}16--24{]}  & \checkmark\ OK \\
        Mexicable        & 20.0 & 17.4 & {[}16--24{]}  & \checkmark\ OK \\
        Metrobús         & 16.3 & 14.1 & {[}13--20{]}  & \checkmark\ OK \\
        Mexibús          & 16.3 & 14.1 & {[}13--20{]}  & \checkmark\ OK \\
        \textsc{rtp}     & 14.0 &  8.0 & {[}10--18{]}  & $\approx$ Revisar \\
        Corredores Conc. & 11.0 &  8.0 & {[}8--14{]}   & \checkmark\ OK \\
        Pumabús          & 11.0 &  8.0 & {[}8--14{]}   & \checkmark\ OK \\
        \hline
    \end{tabular}
    \fuentefigura{Elaboración propia con \texttt{VFTModel}
    \citep{cabrera2026VFTModel}. Fuentes de referencia: STC Metro \cdmx{},
    Metrobús \cdmx{}, \textsc{rtp} \cdmx{}, \textsc{ste} Trolebús.}
\end{table}

\begin{figure}[H]
    \centering
    \includegraphics[width=\textwidth]{Figures/Cap5/fig01_impedancia.png}
    \caption[Caracterización de la impedancia temporal por sistema]{(a)~Velocidad
    efectiva post-fricción por sistema, coloreada según el tipo de derecho de
    vía. (b)~Distribución de tiempos de viaje por segmento interestación
    (segmentos menores a 5~min, sin valores atípicos). Corrida: 2026-04-18.}
    \label{fig:impedancia-sistemas}
    \fuentefigura{Elaboración propia con \texttt{VFTModel}
    \citep{cabrera2026VFTModel}.}
\end{figure}

Los resultados son coherentes con la metodología adoptada. El caso más sensible
es el del \textbf{Trolebús}, cuya velocidad efectiva ($\approx$13.0~km/h con
$CF = 1.380$) se sitúa ligeramente por debajo del límite inferior del rango
operativo documentado (14--22~km/h). Esta diferencia de aproximadamente 1~km/h
(7\,\%) es atribuible a que el tipo de derecho de vía \termino{compartido} no
captura la preferencia operativa parcial que el trolebús mantiene en ciertas
vialidades de la \cdmx{}. Para fines de análisis estructural de red, esta
desviación se encuentra dentro de los márgenes de calibración aceptables.

% !TeX root = ../../tesis.tex
%%%%%%%%%%%%%%%%%%%%%%%%%%%%%%%%%%%%%%%%%%%%%%%%%%%%%%%%%%%%%%%%%%%
% SECCIÓN 5.3.2 — NIVEL DE COBERTURA ESPACIAL (C)
%%%%%%%%%%%%%%%%%%%%%%%%%%%%%%%%%%%%%%%%%%%%%%%%%%%%%%%%%%%%%%%%%%%
\subsection{Nivel de Cobertura Espacial (\texorpdfstring{$C$}{C})}
\label{subsec:cobertura-espacial}

El \textbf{Nivel de Cobertura Espacial} ($C$) evalúa la equidad territorial
de la red al cuantificar qué proporción del área de cada demarcación
político-administrativa se encuentra dentro de un radio de caminata aceptable
respecto a alguna estación. Se define como:

\begin{equation}
    C = \frac{A_{\text{cubierta}}}{A_{\text{total}}} \times 100
    \label{eq:cobertura-res}
\end{equation}

donde $A_{\text{cubierta}}$ es el área resultante de la unión de los
isócronos peatonales (\termino{buffers} de 800~m) de todas las estaciones del
sistema sobre la demarcación, y $A_{\text{total}}$ es la superficie geodésica
total de la demarcación analizada. El análisis abarca las 16~alcaldías de la
\cdmx{} y los 125 municipios de la \zmvm{} incluidos en el dataset
del \textbf{Apímetro} \citep{cabrera2025apimetro}, para un total de
141 demarcaciones.

La Figura~\ref{fig:grafica-cobertura-sistema} sintetiza el nivel de cobertura
agregado por sistema, evidenciando la marcada diferencia entre la red capilar
de superficie y los sistemas estructurados de mayor escala.

\begin{figure}[H]
    \centering
    \includegraphics[width=0.92\textwidth]{%
        Figures/Mapas/Tableau_Resources/RedActual/VFTModel/UniGrafica/%
        Grafica_VFTModel_2026_Cobertura_sistema}
    \caption[Cobertura territorial agregada por sistema de transporte]{
    Cobertura territorial agregada por sistema de transporte en la \zmvm{},
    calculada a partir de isócronos peatonales de 800~m sobre las
    141 demarcaciones analizadas. Corrida: 2026-04-18.}
    \label{fig:grafica-cobertura-sistema}
    \fuentefigura{Elaboración propia con \textbf{Tableau};
    datos: \texttt{VFTModel} \citep{cabrera2026VFTModel}.}
\end{figure}

% ---------------------------------------------------------------
\subsubsection{Sistemas de Infraestructura Estructurada}
\label{subsubsec:cobertura-estructurados}
% ---------------------------------------------------------------

Los sistemas con derecho de vía exclusivo o confinado —\textsc{Metro},
Metrobús, Tren Ligero, Cablebús, Tren Suburbano, Mexicable y Mexibús— presentan
coberturas territoriales altas en las alcaldías centrales, pero decrecen
abruptamente hacia la periferia. Su distribución refleja la lógica de inversión
en infraestructura pesada, concentrada históricamente en las zonas de mayor
densidad y actividad económica de la cuenca.

\begin{table}[H]
    \centering
    \caption[Cobertura territorial — sistemas estructurados, top 3 por alcaldía]{
    Tres alcaldías de la \cdmx{} con mayor cobertura territorial por sistema de
    infraestructura estructurada. Isócrono: 800~m. Corrida: 2026-04-18.}
    \label{tab:cobertura-estructurados}
    \fontsize{8.5pt}{10.5pt}\selectfont
    \begin{tabularx}{\textwidth}{|l|X|r|r|r|}
        \hline
        \textbf{Sistema} & \textbf{Alcaldía} &
        \textbf{Área total (km²)} & \textbf{Área cub. (km²)} &
        \textbf{$C$ (\%)} \\
        \hline
        \multirow{3}{*}{\textsc{Metro}}
            & Cuauhtémoc          & 32.50 & 28.35 & 87.23 \\
            & Benito Juárez       & 26.68 & 21.34 & 79.99 \\
            & Venustiano Carranza & 33.84 & 24.39 & 72.77 \\
        \hline
        \multirow{3}{*}{Metrobús}
            & Cuauhtémoc    & 32.50 & 26.74 & 82.27 \\
            & Benito Juárez & 26.68 & 16.21 & 60.74 \\
            & Iztacalco     & 23.08 & 13.95 & 60.44 \\
        \hline
        \multirow{3}{*}{Tren Ligero}
            & Coyoacán   &  53.88 &  9.31 & 17.28 \\
            & Xochimilco & 114.03 &  6.82 &  5.98 \\
            & Tlalpan    & 314.25 &  4.07 &  1.30 \\
        \hline
        \multirow{3}{*}{Cablebús}
            & Gustavo A. Madero &  87.84 & 11.31 & 12.88 \\
            & Miguel Hidalgo    &  46.39 &  5.50 & 11.85 \\
            & Iztapalapa        & 113.07 & 12.74 & 11.27 \\
        \hline
        \multirow{3}{*}{Tren Suburbano}
            & Cuauhtémoc   & 32.50 & 2.01 & 6.18 \\
            & Azcapotzalco & 33.50 & 2.01 & 5.99 \\
            & (resto)      &  ---  & 0.00 & 0.00 \\
        \hline
        \multirow{3}{*}{Mexicable}
            & Gustavo A. Madero &  87.84 & 2.80 & 3.19 \\
            & (resto)           &   ---  & 0.00 & 0.00 \\
            &                   &        &      &      \\
        \hline
        \multirow{3}{*}{Mexibús\textsuperscript{†}}
            & Iztacalco           & 23.08 & 1.52 & 6.58 \\
            & Venustiano Carranza & 33.84 & 1.90 & 5.62 \\
            & Gustavo A. Madero   & 87.84 & 3.23 & 3.67 \\
        \hline
    \end{tabularx}
    \vspace{0.4em}

    \noindent\textsuperscript{†}\,Datos de geometría y paradas integrados
    desde \termino{OpenStreetMap} \citep{osm2012} por ausencia de registro
    en el Portal de Datos Abiertos a la fecha de corrida.

    \fuentefigura{Elaboración propia con \texttt{VFTModel}
    \citep{cabrera2026VFTModel}; datos: \citep{adip_gtfs_2024, osm2012}.}
\end{table}

% ---------------------------------------------------------------
\subsubsection{Transporte de Superficie}
\label{subsubsec:cobertura-superficie}
% ---------------------------------------------------------------

Los sistemas que operan en vía compartida o mixta —\textsc{rtp}, Corredores
Concesionados, Trolebús, Pumabús y Tren Interurbano— presentan el patrón
territorial inverso: coberturas elevadas en las alcaldías centrales, pero con
una capilaridad que penetra en zonas periféricas donde la infraestructura
pesada no llega. La \textsc{rtp} y los Corredores Concesionados son los
sistemas con mayor cobertura absoluta de la red metropolitana, superando el
90\,\% en múltiples alcaldías.

\begin{table}[H]
    \centering
    \caption[Cobertura territorial — transporte de superficie, top 3 por alcaldía]{
    Tres alcaldías de la \cdmx{} con mayor cobertura territorial por sistema
    de transporte de superficie. Isócrono: 800~m. Corrida: 2026-04-18.}
    \label{tab:cobertura-superficie}
    \fontsize{8.5pt}{10.5pt}\selectfont
    \begin{tabularx}{\textwidth}{|l|X|r|r|r|}
        \hline
        \textbf{Sistema} & \textbf{Alcaldía} &
        \textbf{Área total (km²)} & \textbf{Área cub. (km²)} &
        \textbf{$C$ (\%)} \\
        \hline
        \multirow{3}{*}{\textsc{rtp}}
            & Miguel Hidalgo    & 46.39 & 41.90 & 90.32 \\
            & Cuauhtémoc        & 32.50 & 29.32 & 90.22 \\
            & Azcapotzalco      & 33.50 & 30.08 & 89.81 \\
        \hline
        \multirow{3}{*}{Corredores Conc.}
            & Cuauhtémoc          & 32.50 & 32.47 & 99.90 \\
            & Benito Juárez       & 26.68 & 26.32 & 98.65 \\
            & Iztacalco           & 23.08 & 22.25 & 96.42 \\
        \hline
        \multirow{3}{*}{Trolebús\textsuperscript{†}}
            & Cuauhtémoc    & 32.50 & 22.72 & 69.90 \\
            & Benito Juárez & 26.68 & 16.50 & 61.85 \\
            & Azcapotzalco  & 33.50 & 17.48 & 52.18 \\
        \hline
        \multirow{3}{*}{Pumabús}
            & Coyoacán      &  53.88 & 10.85 & 20.13 \\
            & Álvaro Obregón &  95.82 &  2.25 &  2.34 \\
            & Tlalpan        & 314.25 &  0.25 &  0.08 \\
        \hline
        \multirow{3}{*}{Tren Interurbano}
            & Cuajimalpa de Morelos &  71.10 & 1.59 & 2.23 \\
            & Álvaro Obregón        &  95.82 & 0.42 & 0.44 \\
            & (resto)               &   ---  & 0.00 & 0.00 \\
        \hline
    \end{tabularx}
    \vspace{0.4em}

    \noindent\textsuperscript{†}\,La Línea 13 del Trolebús no contaba con
    registro en el Portal de Datos Abiertos a la fecha de corrida; su
    geometría fue integrada desde \termino{OpenStreetMap}
    \citep{osm2012}. Los valores de cobertura excluyen esa línea.

    \fuentefigura{Elaboración propia con \texttt{VFTModel}
    \citep{cabrera2026VFTModel}; datos: \citep{adip_gtfs_2024, osm2012}.}
\end{table}

% ---------------------------------------------------------------
\subsubsection{Zonas de Mayor Déficit Territorial}
\label{subsubsec:cobertura-deficit}
% ---------------------------------------------------------------

El análisis de cobertura integrada —considerando la unión de los isócronos de
todos los sistemas— revela una inequidad territorial estructural: las
demarcaciones con menor acceso al transporte masivo son, de manera consistente,
las de mayor extensión territorial y menor densidad urbana consolidada, ubicadas
en las periferias sur y oriente de la cuenca.

\begin{table}[H]
    \centering
    \caption[Demarcaciones con menor cobertura de transporte masivo]{
    Cinco demarcaciones de la \cdmx{} y del Estado de México con menor
    cobertura territorial integrada del sistema multimodal. Isócrono: 800~m.
    Corrida: 2026-04-18.}
    \label{tab:cobertura-deficit}
    \begin{tabular}{|l|l|r|r|r|}
        \hline
        \textbf{Entidad} & \textbf{Demarcación} &
        \textbf{Área total (km²)} & \textbf{Área cub. (km²)} &
        \textbf{$C$ (\%)} \\
        \hline
        \multirow{5}{*}{\cdmx{}}
            & Milpa Alta              & 298.01 &  35.00 & 11.75 \\
            & La Magdalena Contreras  &  63.37 &  16.33 & 25.76 \\
            & Tlalpan                 & 314.25 &  94.31 & 30.01 \\
            & Tláhuac                 &  85.78 &  43.85 & 51.12 \\
            & Cuajimalpa de Morelos   &  71.10 &  37.12 & 52.20 \\
        \hline
        \multirow{5}{*}{Estado de México}
            & Tonanitla    &   8.52 & 0.00 & 0.00 \\
            & Capulhuac    &  21.50 & 0.00 & 0.00 \\
            & Calimaya     & 104.26 & 0.00 & 0.00 \\
            & Donato Guerra & 181.36 & 0.00 & 0.00 \\
            & Ayapango     &  50.83 & 0.00 & 0.00 \\
        \hline
    \end{tabular}
    \fuentefigura{Elaboración propia con \texttt{VFTModel}
    \citep{cabrera2026VFTModel}.}
\end{table}

La Figura~\ref{fig:grafica-cobertura-alcaldia} complementa esta lectura con
la distribución de cobertura integrada por alcaldía, permitiendo identificar
el gradiente centro-periferia que caracteriza la inequidad territorial de la
red metropolitana.

\begin{figure}[H]
    \centering
    \includegraphics[width=0.92\textwidth]{%
        Figures/Mapas/Tableau_Resources/RedActual/VFTModel/UniGrafica/%
        Grafica_VFTModel_2026_Cobertura_alcaldia}
    \caption[Cobertura territorial integrada por alcaldía de la \cdmx{}]{
    Cobertura territorial integrada de la red multimodal por alcaldía de la
    \cdmx{}, expresada como porcentaje del área cubierta por isócronos
    peatonales de 800~m. El gradiente centro-periferia es el patrón
    dominante.}
    \label{fig:grafica-cobertura-alcaldia}
    \fuentefigura{Elaboración propia con \textbf{Tableau};
    datos: \texttt{VFTModel} \citep{cabrera2026VFTModel}.}
\end{figure}

% ---------------------------------------------------------------
\subsubsection{Vista Cartográfica de la Cobertura Total}
\label{subsubsec:cobertura-cartografia}
% ---------------------------------------------------------------

Las Figuras~\ref{fig:mapa-cobertura-red-total} y~\ref{fig:mapa-cobertura-calor}
sintetizan la distribución espacial de la cobertura integrada mediante dos
representaciones cartográficas complementarias: el alcance territorial acumulado
de los isócronos de todos los sistemas y la densidad de esa cobertura como mapa
de calor. Ambas fueron elaboradas con \textbf{QGIS} a partir de los
\termino{GeoPackages} exportados por el \textbf{Modelo VFT}.

\begin{figure}[H]
    \centering
    \includegraphics[width=0.90\textwidth]{%
        Figures/Mapas/Qgis_Resources/RedActual/Mapa_Cobertura_Red_Total}
    \caption[Cobertura espacial total de la red multimodal — isócronas 800~m]{
    Cobertura espacial integrada de los once sistemas de transporte masivo de
    la \zmvm{}, calculada a partir de isócronos peatonales de 800~m. Las zonas
    sin cobertura se concentran en el sur y oriente de la cuenca metropolitana,
    evidenciando las demarcaciones con mayor exclusión territorial.}
    \label{fig:mapa-cobertura-red-total}
    \fuentefigura{Elaboración propia con \textbf{QGIS};
    datos: \texttt{VFTModel} \citep{cabrera2026VFTModel},
    GTFS \citep{adip_gtfs_2024} e INEGI 2023.}
\end{figure}

\begin{figure}[H]
    \centering
    \includegraphics[width=0.90\textwidth]{%
        Figures/Mapas/Qgis_Resources/RedActual/Mapa_Cobertura_Mapa_Calor}
    \caption[Densidad de cobertura de la red multimodal — mapa de calor]{
    Densidad de la cobertura espacial de la red multimodal expresada como mapa
    de calor. Los núcleos de máxima accesibilidad se ubican en las alcaldías
    centrales (Cuauhtémoc, Benito Juárez, Iztacalco), mientras que las
    periferias sur y oriente presentan densidades cercanas a cero, confirmando
    la inequidad territorial diagnosticada.}
    \label{fig:mapa-cobertura-calor}
    \fuentefigura{Elaboración propia con \textbf{QGIS};
    datos: \texttt{VFTModel} \citep{cabrera2026VFTModel},
    GTFS \citep{adip_gtfs_2024} e INEGI 2023.}
\end{figure}

Los datos de los once sistemas apuntan en la misma dirección: el acceso al
transporte masivo en la \zmvm{} no se distribuye de manera homogénea. Alcaldías
como Cuauhtémoc registran coberturas de hasta el 99.90\,\% bajo los Corredores
Concesionados; Milpa Alta, segunda demarcación en superficie de la \cdmx{}
solo superada por Tlalpan, no rebasa el 11.75\,\% en ningún sistema. Esa
distancia no expresa una diferencia operativa sino una acumulación histórica de
trazados e inversiones que reproduce, dentro del campo de la movilidad urbana,
la estratificación descrita por \citet{rendon2009ciudadanos}: ciudadanos con
acceso pleno a la red de oportunidades de la ciudad frente a quienes lo tienen
de forma parcial o ninguna. Sobre esta base territorial, la
Sección~\ref{subsec:fuerza-capilar} examina cómo esa escasez periférica se
concentra en un puñado de nodos que conectan la periferia con la red troncal.

% !TeX root = ../../tesis.tex
%%%%%%%%%%%%%%%%%%%%%%%%%%%%%%%%%%%%%%%%%%%%%%%%%%%%%%%%%%%%%%%%%%%
% SECCIÓN 5.3.3 — FUERZA CAPILAR (FC)
%%%%%%%%%%%%%%%%%%%%%%%%%%%%%%%%%%%%%%%%%%%%%%%%%%%%%%%%%%%%%%%%%%%
\subsection{Fuerza Capilar (\texorpdfstring{$FC$}{FC})}
\label{subsec:fuerza-capilar}

Todo usuario del transporte metropolitano reconoce el fenómeno
aunque no lo nombre: el nodo de Barranca del Muerto, el complejo de
Chapultepec, los grandes paraderos sobre Tlalpan. Son puntos donde
múltiples líneas y sistemas convergen, donde la red se aprieta antes
de abrirse de nuevo hacia las colonias periféricas. La \textbf{Fuerza
Capilar} ($FC$) es el indicador que formaliza esa intuición urbana:
cuantifica el nivel de convergencia de flujos en cada nodo o
agrupación de nodos del grafo, identificando los puntos donde la red
concentra y redistribuye la mayor cantidad de conexiones. El
indicador se evalúa en dos perspectivas complementarias: la
\textbf{Fuerza Capilar de Macro-Hub} ($FC_H$), que agrupa estaciones
físicamente colindantes en un radio de 100~m para representar los
grandes nodos de transbordo, y la \textbf{Fuerza Capilar individual}
($FC$), que mide la convergencia topológica estricta en cada nodo sin
tolerancia espacial.

La expresión formal de la perspectiva de Macro-Hub es:

\begin{equation}
    FC_{H} = \sum_{j \in H}
    \bigl(k_{\text{in}}(j) + k_{\text{out}}(j)\bigr)
    - \text{Aristas Internas}
    \label{eq:fc-geografica}
\end{equation}

\noindent donde:
\begin{itemize}
    \item $H$: conjunto de nodos del Macro-Hub (estaciones agrupadas
          en un radio de tolerancia de 100~m).
    \item $k_{\text{in}}(j)$: número de aristas entrantes del nodo $j$
          dentro del grafo de servicios.
    \item $k_{\text{out}}(j)$: número de aristas salientes del nodo $j$.
    \item Aristas Internas: conexiones entre nodos del propio Hub,
          descontadas para evitar el doble conteo de flujos internos
          al complejo.
\end{itemize}

% ---------------------------------------------------------------
\subsubsection{Distribución por Nodo y Macro-Hub}
\label{subsubsec:fc-distribucion-nodos}
% ---------------------------------------------------------------

La Figura~\ref{fig:grafica-fc-nodos} sintetiza la distribución de la
Fuerza Capilar en los principales nodos de la red, ordenados por su
nivel de concentración de conexiones.

\begin{figure}[H]
    \centering
    \includegraphics[width=0.92\textwidth]{%
        Figures/Mapas/Tableau_Resources/RedActual/VFTModel/UniGrafica/%
        Grafica_VFTModel_2026_Fuerza_capilar}
    \caption[Distribución de Fuerza Capilar en los principales nodos
    de la red]{%
    Distribución de la Fuerza Capilar ($FC$ y $FC_H$) en los
    principales nodos y Macro-Hubs de transferencia de la \zmvm{}.
    Los valores más altos corresponden a los complejos donde mayor
    número de rutas y sistemas se superponen. Corrida: 2026-04-18.}
    \label{fig:grafica-fc-nodos}
    \fuentefigura{Elaboración propia con \textbf{Tableau};
    datos: \texttt{VFTModel} \citep{cabrera2026VFTModel}.}
\end{figure}

% ---------------------------------------------------------------
\subsubsection{Distribución Territorial}
\label{subsubsec:fc-distribucion-territorial}
% ---------------------------------------------------------------

Leída en conjunto con el análisis de cobertura de la sección
anterior, la distribución de la Fuerza Capilar revela la otra cara
del déficit territorial: las demarcaciones con menor cobertura no
carecen simplemente de estaciones, sino de nodos redistributores
capaces de articular los ramales locales con la red troncal. Las
zonas periféricas que registran los menores valores de $C$ son
también las que presentan menor densidad de nodos de alta $FC$: en
ellas, la red no solo es escasa, sino que los pocos nodos existentes
cargan con una proporción desproporcionada de las conexiones locales,
sin el respaldo de un Macro-Hub que diversifique la oferta de
transferencia.

La Figura~\ref{fig:mapa-fuerza-capilar} espacializa esa distribución
sobre la cuenca metropolitana, permitiendo identificar los corredores
y zonas de mayor densidad capilar a escala territorial.

\begin{figure}[H]
    \centering
    \includegraphics[width=0.90\textwidth]{%
        Figures/Mapas/Qgis_Resources/RedActual/Mapa_Fuerza_Capilar}
    \caption[Distribución espacial de la Fuerza Capilar — red
    multimodal \zmvm{}]{%
    Distribución espacial de la Fuerza Capilar ($FC_H$) en la red
    de transporte masivo de la \zmvm{}. Los nodos de mayor
    concentración de conexiones se ubican en el corredor sur de la
    cuenca, sobre los ejes Insurgentes--Tlalpan y el entorno del
    Anillo Periférico. Corrida: 2026-04-18.}
    \label{fig:mapa-fuerza-capilar}
    \fuentefigura{Elaboración propia con \textbf{QGIS};
    datos: \texttt{VFTModel} \citep{cabrera2026VFTModel},
    GTFS \citep{adip_gtfs_2024} e INEGI 2023.}
\end{figure}

% ---------------------------------------------------------------
\subsubsection{Vista Integral de la Red}
\label{subsubsec:fc-vista-integral}
% ---------------------------------------------------------------

La Figura~\ref{fig:dashboard-fuerza-capilar} consolida, mediante el
\termino{dashboard} analítico del modelo, la visión integral del
indicador sobre el conjunto de la red.

\begin{figure}[H]
    \centering
    \includegraphics[width=0.95\textwidth]{%
        Figures/Mapas/Tableau_Resources/RedActual/VFTModel/Dashboards/%
        Dashboard_VFTModel_2026_Fuerza_Capilar}
    \caption[Dashboard de Fuerza Capilar — red multimodal \zmvm{}]{%
    Vista integral de la Fuerza Capilar ($FC$ y $FC_H$) sobre la red
    de transporte masivo de la \zmvm{}, generada con el
    \termino{dashboard} analítico del \texttt{VFTModel}.
    Corrida: 2026-04-18.}
    \label{fig:dashboard-fuerza-capilar}
    \fuentefigura{Elaboración propia con \textbf{Tableau};
    datos: \texttt{VFTModel} \citep{cabrera2026VFTModel}.}
\end{figure}

Los nodos de alta Fuerza Capilar trazan, sobre la ciudad, una
jerarquía implícita que no emergió de una planificación orientada al
transporte colectivo: en una metrópoli cuya infraestructura vial se
estructuró históricamente en función del automóvil privado, la red de
transporte masivo se fue articulando sobre los márgenes e intersticios
de esa trama, y su capacidad redistributiva terminó concentrada en un
puñado de complejos del corredor sur de los que dependen decenas de
rutas periféricas. Esa concentración no es neutral; reproduce, en
términos de conectividad, la misma asimetría territorial documentada
en la cobertura. La sección siguiente evalúa su consecuencia directa
sobre el usuario: el Factor de Desviación mide cuánto se alarga el
viaje cuando esa jerarquía obliga a rodeos.

% !TeX root = ../../tesis.tex
%%%%%%%%%%%%%%%%%%%%%%%%%%%%%%%%%%%%%%%%%%%%%%%%%%%%%%%%%%%%%%%%%%%
% SECCIÓN 5.3.4 — FACTOR DE DESVIACIÓN (DI)
%%%%%%%%%%%%%%%%%%%%%%%%%%%%%%%%%%%%%%%%%%%%%%%%%%%%%%%%%%%%%%%%%%%
\subsection{Factor de Desviación (\texorpdfstring{$DI$}{DI})}
\label{subsec:factor-desviacion}

Cuando alguien planea un viaje con el mapa en mano, la distancia en
línea recta entre origen y destino parece razonable. Lo que la red de
transporte impone en la práctica es otra cosa: rutas que obligan a
converger hacia el centro antes de salir hacia la periferia opuesta,
trayectos que multiplican el recorrido respecto al camino ideal. El
\textbf{Factor de Desviación} ($DI$) cuantifica exactamente esa
brecha: es la razón entre la distancia real recorrida sobre la
infraestructura —incluyendo la primera y última milla peatonal— y la
distancia teórica en línea recta entre el mismo par origen-destino.

\begin{equation}
    DI = \frac{D_{\text{red}} + D_{\text{caminata}}}
              {D_{\text{haversine}}(O, D)}
    \label{eq:factor-desviacion}
\end{equation}

\noindent donde:
\begin{itemize}
    \item $D_{\text{red}}$: distancia acumulada sobre el trazo real de
          la ruta óptima en tiempo.
    \item $D_{\text{caminata}}$: suma de la primera y última milla
          peatonal (acceso y egreso a la red).
    \item $D_{\text{haversine}}(O, D)$: distancia en línea recta entre
          origen $O$ y destino $D$, calculada con la fórmula de
          Haversine.
\end{itemize}

Un valor de $DI = 1.0$ indica una ruta perfectamente directa; valores
superiores a $1.5$ revelan ineficiencias topológicas significativas en
el par origen-destino analizado.

% ---------------------------------------------------------------
\subsubsection{Distribución en la Red}
\label{subsubsec:di-distribucion}
% ---------------------------------------------------------------

Sobre una muestra de \textbf{100 pares origen-destino} (semilla
reproducible $= 42$), el modelo arrojó los estadísticos globales del
Cuadro~\ref{tab:di-estadisticas}. La distribución resultante —con
media de $1.545$ y valores extremos entre $0.97$ y $4.34$— confirma
que el recorrido promedio en la red metropolitana supera en más del
50\,\% la distancia ideal.

\begin{table}[H]
    \centering
    \caption[Estadísticos globales del Factor de Desviación]{%
    Estadísticos descriptivos del $DI$ sobre 100 pares
    origen-destino de la red de la \zmvm{}.
    Corrida: 2026-04-27.}
    \label{tab:di-estadisticas}
    \begin{tabular}{|l|r|}
        \hline
        \textbf{Estadístico} & \textbf{Valor} \\
        \hline
        $n$ (rutas calculadas)     & 100   \\
        Media ($\bar{DI}$)         & 1.545 \\
        Desviación estándar        & 0.517 \\
        Mínimo                     & 0.97  \\
        Percentil 25               & 1.27  \\
        Mediana                    & 1.40  \\
        Percentil 75               & 1.67  \\
        Máximo                     & 4.34  \\
        \hline
    \end{tabular}
    \fuentefigura{Elaboración propia con \texttt{VFTModel}
    \citep{cabrera2026VFTModel}. Corrida: 2026-04-27.}
\end{table}

El registro completo de los 100 pares analizados —con sus distancias
individuales, valores de $DI$ y número de saltos— se incluye en el
Cuadro~\ref{tab:anx-di-corrida-completa} del
Anexo~\ref{anx:visualizaciones-analiticas}.

La Figura~\ref{fig:mapa-di} sitúa esa distribución sobre el
territorio, permitiendo leer la magnitud del rodeo como fenómeno
espacial antes de desagregarla por alcaldía.

\begin{figure}[H]
    \centering
    \includegraphics[width=0.90\textwidth]{%
        Figures/Mapas/Qgis_Resources/RedActual/Mapa_DF}
    \caption[Distribución espacial del Factor de Desviación ($DI$)
    por nodo]{%
    Factor de Desviación ($DI$) por nodo en la red de transporte
    masivo de la \zmvm{}. El gradiente verde-a-rojo señala las
    zonas donde la infraestructura impone mayores rodeos respecto
    a la distancia lineal ideal. Corrida: 2026-04-27.}
    \label{fig:mapa-di}
    \fuentefigura{Elaboración propia con \textbf{QGIS};
    datos: \texttt{VFTModel} \citep{cabrera2026VFTModel},
    GTFS \citep{adip_gtfs_2024} e INEGI 2023.}
\end{figure}

La Figura~\ref{fig:grafica-di-alcaldia} desagrega esa imagen
territorial mostrando el $DI$ promedio por alcaldía, evidenciando
cuáles demarcaciones concentran la mayor ineficiencia de trayecto.

\begin{figure}[H]
    \centering
    \includegraphics[width=0.92\textwidth]{%
        Figures/Mapas/Tableau_Resources/RedActual/VFTModel/UniGrafica/%
        Grafica_VFTModel_2026_Factor_desviacion}
    \caption[Factor de Desviación promedio por alcaldía de la
    \zmvm{}]{%
    Factor de Desviación ($DI$) promedio por alcaldía de la
    \zmvm{}, calculado sobre la muestra de 100 pares
    origen-destino. Las alcaldías periféricas del oriente y sur
    presentan los valores más altos, reflejando la ausencia de
    conectividad transversal en esos cuadrantes.
    Corrida: 2026-04-27.}
    \label{fig:grafica-di-alcaldia}
    \fuentefigura{Elaboración propia con \textbf{Tableau};
    datos: \texttt{VFTModel} \citep{cabrera2026VFTModel}.}
\end{figure}

El patrón es consistente con los indicadores anteriores: las
demarcaciones que registran menor cobertura y menor densidad de
nodos de alta Fuerza Capilar son también las que imponen los mayores
rodeos a sus usuarios.

% ---------------------------------------------------------------
\subsubsection{Casos de Estudio}
\label{subsubsec:di-casos}
% ---------------------------------------------------------------

Para anclar los estadísticos en viajes concretos, se seleccionaron
doce pares origen-destino estratégicos concentrados en las periferias
sur y oriente de la cuenca. El Cuadro~\ref{tab:di-casos} consolida
los valores de $DI$ para cada caso.

\begin{table}[H]
    \centering
    \caption[Factor de Desviación — doce casos de estudio]{%
    Factor de Desviación ($DI$) para doce pares origen-destino
    estratégicos de la \zmvm{}, seleccionados en las periferias
    sur y oriente. Corrida: 2026-04-27.}
    \label{tab:di-casos}
    \small
    \begin{tabularx}{\textwidth}{|l|X|X|r|r|r|}
        \hline
        \textbf{\#} & \textbf{Origen} & \textbf{Destino} &
        \textbf{$D_{\text{red}}$ (km)} &
        \textbf{$D_{\text{Hav.}}$ (km)} &
        \textbf{$DI$} \\
        \hline
         1 & Milpa Alta        & Autódromo / Centro-Oriente   & 29.18 & 22.59 & 1.29 \\
         2 & Milpa Alta        & Paseo de la Reforma (CBD)    & 37.73 & 27.50 & 1.37 \\
         3 & Milpa Alta        & Santa Fe (CBD)               & 36.42 & 28.62 & 1.27 \\
         4 & Milpa Alta        & FES Acatlán                  & 48.34 & 35.73 & 1.35 \\
         5 & Milpa Alta        & Ciudad Universitaria         & 26.25 & 19.96 & 1.32 \\
         6 & Milpa Alta        & Coyoacán                     & 25.24 & 18.86 & 1.34 \\
         7 & Santa Catarina    & Ciudad Universitaria         & 42.06 & 23.10 & 1.82 \\
         8 & Santa Catarina    & FES Acatlán                  & 45.64 & 32.16 & 1.42 \\
         9 & Milpa Alta        & Santa Catarina (Sur-Oriente) & 37.22 & 14.79 & 2.52 \\
        10 & Milpa Alta        & Ajusco                       & 29.19 & 20.56 & 1.42 \\
        11 & Ajusco            & Preparatoria 1               & 15.91 & 11.16 & 1.43 \\
        12 & Parque Los Olivos & Centro de Tláhuac            &  2.20 &  2.08 & 1.06 \\
        \hline
        \multicolumn{5}{|l|}{\textbf{Mejor caso (muestra aleatoria):} Plutarco Elías Calles $\to$ Fuente de Loreto} & \textbf{0.97} \\
        \multicolumn{5}{|l|}{\textbf{Peor caso (muestra aleatoria):} Carlos A. Vidal $\to$ 20 de Noviembre} & \textbf{4.34} \\
        \hline
    \end{tabularx}
    \fuentefigura{Elaboración propia con \texttt{VFTModel}
    \citep{cabrera2026VFTModel}. Corrida: 2026-04-27.}
\end{table}

Entre los doce casos, el Caso~9 registra el $DI$ más alto del
conjunto seleccionado ($DI = 2.52$): para ir de Milpa Alta a Santa
Catarina Yecahuizotl —dos puntos del sur-oriente separados por apenas
14.8~km en línea recta— la red obliga a recorrer 37.2~km, más del
doble de la distancia ideal. Ese valor no es un accidente de la
muestra; es la expresión métrica de la ausencia de conectividad
transversal en ese cuadrante de la cuenca.

% ---------------------------------------------------------------
\subsubsection{Vista Integral de la Red}
\label{subsubsec:di-vista-integral}
% ---------------------------------------------------------------

¿Cómo se distribuye el costo del rodeo sobre el territorio? La
Figura~\ref{fig:dashboard-di} responde esa pregunta articulando en
un solo panel la magnitud del $DI$ por alcaldía, la dispersión
estadística de la muestra y la huella geográfica de las rutas
analizadas.

\begin{figure}[H]
    \centering
    \includegraphics[width=0.95\textwidth]{%
        Figures/Mapas/Tableau_Resources/RedActual/VFTModel/Dashboards/%
        Dashboard_VFTModel_2026_Factor_desviacion}
    \caption[Dashboard del Factor de Desviación — red multimodal
    \zmvm{}]{%
    Vista integral del Factor de Desviación ($DI$) sobre la red de
    transporte masivo de la \zmvm{}, generada con el
    \termino{dashboard} analítico del \texttt{VFTModel}.
    Corrida: 2026-04-27.}
    \label{fig:dashboard-di}
    \fuentefigura{Elaboración propia con \textbf{Tableau};
    datos: \texttt{VFTModel} \citep{cabrera2026VFTModel}.}
\end{figure}

El análisis del \texttt{VFTModel} identifica en Tláhuac la alcaldía
donde ese costo estructural es más pronunciado. El Caso~7 lo
cuantifica con precisión: de Santa Catarina Yecahuizotl a Ciudad
Universitaria hay 23~km en línea recta, pero la red exige 42~km de
recorrido real ($DI = 1.82$), porque ninguna ruta de los datos
procesados por \texttt{Apímetro} conecta esos dos puntos del sur sin
internar el trayecto hacia el centro de la ciudad. No es un caso
límite ni un atípico estadístico: es la geometría de una red sin
conectividad transversal en ese cuadrante, medida con los mismos
instrumentos que describen el resto del capítulo.

\begin{figure}[H]
    \centering
    \includegraphics[width=0.85\textwidth]{%
        Figures/Cap5/FactorDesviacion/Mapa/NB04_mapa10_caso07.png}
    \caption[Caso~7: Santa Catarina Yecahuizotl → Ciudad
    Universitaria ($DI = 1.82$)]{%
    Caso~7: ruta óptima de Santa Catarina Yecahuizotl a Ciudad
    Universitaria. La red obliga a un recorrido de 42.06~km para
    cubrir una distancia lineal de 23.10~km ($DI = 1.82$), al no
    existir conectividad transversal directa en el cuadrante
    sur-oriente de la \zmvm{}. Corrida: 2026-04-27.}
    \label{fig:di-caso07}
    \fuentefigura{Elaboración propia con \texttt{VFTModel}
    \citep{cabrera2026VFTModel}.}
\end{figure}

Los cuatro indicadores de Fase~1 —impedancia, cobertura, Fuerza
Capilar y Factor de Desviación— describen colectivamente una red
metropolitana con capacidades desiguales: bien articulada en el
centro, con déficit progresivo hacia la periferia sur y oriente, y
sin conectividad transversal en ese cuadrante. Esa descripción es
topológica: se basa en la estructura del grafo, sin considerar las
condiciones reales de operación. La Fase~2 incorpora esas condiciones
mediante dos parámetros de ponderación dinámica —el Coeficiente de
Fricción Vial ($CF$) y la penalización por transferencia ($T_b$)—
que ajustan el modelo a la experiencia de viaje efectiva.


% !TeX root = ../../tesis.tex
%%%%%%%%%%%%%%%%%%%%%%%%%%%%%%%%%%%%%%%%%%%%%%%%%%%%%%%%%%%%%%%%%%%
% SECCIÓN 5.4 — PONDERACIÓN DINÁMICA: FASE 2
%%%%%%%%%%%%%%%%%%%%%%%%%%%%%%%%%%%%%%%%%%%%%%%%%%%%%%%%%%%%%%%%%%%
\section{Ponderación Dinámica: Fase~2}
\label{sec:fase2-ponderacion}

Los indicadores de Fase~1 caracterizaron la red desde su geometría y
conectividad internas. Para aproximar los tiempos de viaje reales, el
modelo incorpora dos parámetros que ajustan los pesos del grafo a las
condiciones operativas observadas: el Coeficiente de Fricción Vial
($CF$), que diferencia la velocidad efectiva según el derecho de vía de
cada sistema, y la penalización por transferencia ($T_b$), que
contabiliza el costo temporal de cada trasbordo. Juntos, estos
parámetros preparan el grafo ponderado sobre el que la Fase~3 calcula
los tiempos de viaje y la centralidad de la red.

% !TeX root = ../../tesis.tex
%%%%%%%%%%%%%%%%%%%%%%%%%%%%%%%%%%%%%%%%%%%%%%%%%%%%%%%%%%%%%%%%%%%
% SECCIÓN 5.4.1 — COEFICIENTE DE FRICCIÓN VIAL (CF)
%%%%%%%%%%%%%%%%%%%%%%%%%%%%%%%%%%%%%%%%%%%%%%%%%%%%%%%%%%%%%%%%%%%
\subsection{Coeficiente de Fricción Vial (\texorpdfstring{$CF$}{CF})}
\label{subsec:fase2-cf}

El Coeficiente de Fricción Vial, definido en §\ref{subsec:friccion-vial}
e implementado como \texttt{FrictionCalculator} en §\ref{subsec:vft-cf},
traduce el tipo de derecho de vía de cada segmento en un multiplicador
sobre el tiempo de traslado:

\begin{equation*}
    CF = 1 + \alpha \cdot \beta_{\text{ZMVM}}
    \tag{\ref{eq:vft-cf}}
\end{equation*}

Con $\beta_{\text{ZMVM}} = 0.759$ \citep{tomtom2024} y cuatro valores
de $\alpha$ calibrados según el grado de segregación física de la
infraestructura \citep{sedatu2019manual}, los coeficientes resultantes
para cada categoría son los del Cuadro~\ref{tab:cf-derecho-via}.

\begin{table}[H]
    \centering
    \caption[Coeficiente de Fricción Vial por categoría de derecho de vía]{%
    Coeficiente de Fricción Vial ($CF$) por categoría de derecho de vía
    y sistemas de transporte asignados en la red de la \zmvm{}.
    Un $CF = 1.000$ indica operación sin interferencia del tráfico.
    Corrida: 2026-04-18.}
    \label{tab:cf-derecho-via}
    \small
    \begin{tabular}{|l|c|r|l|}
        \hline
        \textbf{Derecho de vía} & \textbf{$\alpha$} &
        \textbf{$CF$} & \textbf{Sistemas} \\
        \hline
        Exclusivo  & 0.0 & 1.000 & Metro, Metrobús, Tren Ligero, Tren Suburbano \\
        Confinado  & 0.2 & 1.152 & Cablebús, Mexibús, Mexicable \\
        Compartido & 0.5 & 1.380 & Trolebús \\
        Mixto      & 1.0 & 1.759 & RTP, Corredores Concesionados, Pumabús \\
        \hline
    \end{tabular}
    \fuentefigura{Elaboración propia con \texttt{VFTModel}
    \citep{cabrera2026VFTModel};
    $\beta$: \citep{tomtom2024};
    $\alpha$: \citep{sedatu2019manual}.}
\end{table}

La brecha entre extremos es considerable: los sistemas con derecho de
vía exclusivo operan a su velocidad de diseño sin penalización, mientras
que RTP y los Corredores Concesionados —que concentran más del
85\,\% de las aristas de servicio real de la red— absorben el índice
de congestión completo de la ciudad. Dado que esas líneas son también
las que conectan las periferias con el centro, el $CF$ amplifica las
asimetrías de cobertura y capilaridad ya documentadas en la Fase~1.

% !TeX root = ../../tesis.tex
%%%%%%%%%%%%%%%%%%%%%%%%%%%%%%%%%%%%%%%%%%%%%%%%%%%%%%%%%%%%%%%%%%%
% SECCIÓN 5.4.2 — PENALIZACIÓN POR TRANSFERENCIA (Tb)
%%%%%%%%%%%%%%%%%%%%%%%%%%%%%%%%%%%%%%%%%%%%%%%%%%%%%%%%%%%%%%%%%%%
\subsection{Penalización por Transferencia (\texorpdfstring{$T_b$}{Tb})}
\label{subsec:fase2-tb}

Cambiar de línea o de sistema no es gratuito en términos de tiempo.
Más allá de la caminata entre andenes, el usuario debe esperar la
llegada del siguiente servicio: una espera que depende de la frecuencia
del sistema al que transborda. La penalización por transferencia ($T_b$)
incorpora ese costo al grafo asignando un peso adicional a cada arista
de tipo trasbordo.

\begin{equation}
    T_b = \frac{F}{2}
    \label{eq:penalizacion-transferencia}
\end{equation}

\noindent donde $F$ es la frecuencia del sistema destino expresada en
minutos entre servicios. El factor $\tfrac{1}{2}$ representa la
espera media esperada bajo llegadas aleatorias al andén: en promedio,
el usuario llega en el punto medio del intervalo entre dos servicios
consecutivos.

La penalización actúa exclusivamente sobre las 20,482 aristas de
trasbordo peatonal del grafo —aquellas que conectan nodos de sistemas
distintos dentro de un mismo complejo de transferencia— y no modifica
el tiempo de las aristas de servicio real. Con $T_b$ aplicado, el
grafo refleja no solo el tiempo de traslado sobre cada línea, sino el
costo acumulado de cada cambio de sistema que la ruta óptima imponga
al usuario.

Ambos parámetros —$CF$ y $T_b$— quedan incorporados al grafo ponderado
antes de ejecutar el algoritmo de ruta mínima. La Fase~3 opera sobre
ese grafo para derivar el tiempo de viaje promedio de la red y su
distribución territorial.


% !TeX root = ../../tesis.tex
%%%%%%%%%%%%%%%%%%%%%%%%%%%%%%%%%%%%%%%%%%%%%%%%%%%%%%%%%%%%%%%%%%%
% SECCIÓN 5.5 — ANÁLISIS DE RED: FASE 3
%%%%%%%%%%%%%%%%%%%%%%%%%%%%%%%%%%%%%%%%%%%%%%%%%%%%%%%%%%%%%%%%%%%
\section{Análisis de Red: Fase~3}
\label{sec:fase3-analisis-red}

Las Fases~1 y 2 describen la estructura y los pesos de la red; la
Fase~3 responde una pregunta distinta: ¿cuánto tarda realmente un
viaje en esta red, qué nodos sostienen ese tiempo y qué ocurre cuando
alguno falla? Los tres indicadores de esta fase —tiempo de viaje
promedio, centralidad de intermediación y vulnerabilidad estructural—
operan sobre el componente gigante del grafo (10,561 nodos, 95\,\% de
la red) con el algoritmo de \dijkstra{} ponderado.

% !TeX root = ../../tesis.tex
%%%%%%%%%%%%%%%%%%%%%%%%%%%%%%%%%%%%%%%%%%%%%%%%%%%%%%%%%%%%%%%%%%%
% SECCIÓN 5.5.1 — TIEMPO DE VIAJE PROMEDIO (T)
%%%%%%%%%%%%%%%%%%%%%%%%%%%%%%%%%%%%%%%%%%%%%%%%%%%%%%%%%%%%%%%%%%%
\subsection{Tiempo de Viaje Promedio (\texorpdfstring{$T$}{T})}
\label{subsec:tiempo-viaje-promedio}

Ejecutado el algoritmo de \dijkstra{} ponderado sobre el componente
gigante, el tiempo de viaje promedio de la red resulta en
$T = 108.92$ minutos (corrida: 2026-09-02, 10,561~nodos, tiempo de
cómputo: 251.8~segundos). El valor sintetiza el costo temporal medio
de todos los pares origen-destino posibles bajo las condiciones de
ponderación de Fase~2.

La referencia de campo más cercana proviene de la Encuesta Nacional de
Ocupación y Empleo (ENOE), que estima en aproximadamente 85~minutos el
tiempo promedio de traslado en la \zmvm{}. Frente a los 108.92~minutos
del modelo, la brecha de $\sim$24~minutos tiene origen metodológico, no
en una sobreestimación del sistema. Mientras la ENOE construye su
promedio a partir de desplazamientos observados —concentrados en los
corredores de alta afluencia que conectan el centro con la periferia
próxima— el modelo asigna el mismo peso a cada par del componente
gigante. Los trayectos entre periferias opuestas, que la demanda
cotidiana apenas realiza, quedan contabilizados en igualdad de
condiciones que los recorridos más frecuentes, elevando el promedio
topológico por encima del promedio ponderado por demanda.

Esa variación queda desglosada por alcaldía en la
Figura~\ref{fig:grafica-t-promedio}.

\begin{figure}[H]
    \centering
    \includegraphics[width=0.92\textwidth]{%
        Figures/Mapas/Tableau_Resources/RedActual/VFTModel/UniGrafica/%
        Grafica_VFTModel_2026_Tiempo_Promedio}
    \caption[Tiempo de viaje promedio por alcaldía --- \zmvm{}]{%
    Tiempo de viaje promedio ($T$) por alcaldía. Corrida sobre el
    componente gigante (10,561~nodos) con \dijkstra{} ponderado.
    \zmvm{}, 2026-09-02.}
    \label{fig:grafica-t-promedio}
    \fuentefigura{Elaboración propia con \textbf{Tableau};
    datos: \texttt{VFTModel} \citep{cabrera2026VFTModel}.}
\end{figure}

Con $T$ calculado, surge una pregunta estructural: ¿qué nodos
concentran mayor responsabilidad sobre ese valor? No todos los nodos
son equivalentes en la topología del grafo; algunos articulan una
fracción desproporcionada de las rutas mínimas y su fallo impacta los
tiempos en toda la red. La centralidad de intermediación $B(v)$
—analizada en la sección siguiente— identifica esos puntos.

% !TeX root = ../../tesis.tex
%%%%%%%%%%%%%%%%%%%%%%%%%%%%%%%%%%%%%%%%%%%%%%%%%%%%%%%%%%%%%%%%%%%
% SECCIÓN 5.5.2 — CENTRALIDAD DE INTERMEDIACIÓN B(v)
%%%%%%%%%%%%%%%%%%%%%%%%%%%%%%%%%%%%%%%%%%%%%%%%%%%%%%%%%%%%%%%%%%%
\subsection{Centralidad de Intermediación (\texorpdfstring{$B(v)$}{B(v)})}
\label{subsec:fase3-centralidad}

No todos los nodos de la red tienen el mismo peso en la conectividad
global. Algunos funcionan como puentes inevitables: la mayoría de las
rutas mínimas entre cualquier par de la red pasan por ellos, y su
ausencia fragmentaría la estructura o alargaría los tiempos de forma
generalizada. La centralidad de intermediación $B(v)$ mide exactamente
esa condición: la proporción de rutas mínimas del grafo que atraviesan
cada nodo.

Dado un grafo $\mathcal{G}$ con nodos $\mathcal{V}$, la centralidad de
intermediación de un nodo $v$ se define como:

\begin{equation*}
    B(v) = \sum_{s \neq v \neq t} \frac{\sigma_{st}(v)}{\sigma_{st}}
    \tag{\ref{eq:betweenness}}
\end{equation*}

\noindent donde:
\begin{itemize}
    \item $\sigma_{st}$: cantidad de rutas mínimas entre los nodos $s$ y $t$.
    \item $\sigma_{st}(v)$: subconjunto de esas rutas que pasan por $v$.
\end{itemize}

El modelo calcula $B(v)$ mediante el algoritmo de Brandes, que reduce la
complejidad computacional a $\mathcal{O}(VE)$ sobre grafos ponderados; la
corrida sobre el componente gigante (10,561~nodos) tomó 593.4~segundos.

Dicho de otro modo: hay estaciones que sostienen la red desde adentro
sin que los usuarios lo noten. No son necesariamente las más
transitadas; son las que aparecen en la ruta cuando no hay otra
opción. Quitarlas no molesta a unos cuantos pasajeros —desorganiza
trayectos que atraviesan cuadrantes enteros de la metrópoli.

El Cuadro~\ref{tab:bv-top10} presenta los diez nodos con mayor $B(v)$
calculados sobre el componente gigante. El resultado más llamativo:
los diez son estaciones del Metro, y Tacubaya encabeza la lista con
un valor de 0.2179.

\begin{table}[H]
    \centering
    \caption[Top 10 nodos por centralidad de intermediación $B(v)$]{%
    Diez nodos con mayor centralidad de intermediación ($B(v)$) en la
    red de transporte de la \zmvm{}. Corrida: 2026-09-02.}
    \label{tab:bv-top10}
    \small
    \begin{tabular}{|r|l|l|l|r|}
        \hline
        \textbf{\#} & \textbf{Estación} & \textbf{Sistema} &
        \textbf{Alcaldía} & \textbf{$B(v)$} \\
        \hline
         1 & Tacubaya        & Metro & Miguel Hidalgo & 0.2179 \\
         2 & Mixcoac         & Metro & Benito Juárez  & 0.1725 \\
         3 & Hidalgo         & Metro & Cuauhtémoc     & 0.1595 \\
         4 & Balderas        & Metro & Cuauhtémoc     & 0.1433 \\
         5 & Lázaro Cárdenas & Metro & Cuauhtémoc     & 0.1376 \\
         6 & Chilpancingo    & Metro & Cuauhtémoc     & 0.1366 \\
         7 & Viaducto        & Metro & Benito Juárez  & 0.1346 \\
         8 & Villa de Cortés & Metro & Benito Juárez  & 0.1341 \\
         9 & Nativitas       & Metro & Benito Juárez  & 0.1299 \\
        10 & Zapata          & Metro & Benito Juárez  & 0.1283 \\
        \hline
    \end{tabular}
    \fuentefigura{Elaboración propia con \texttt{VFTModel}
    \citep{cabrera2026VFTModel}. Corrida: 2026-09-02.}
\end{table}

El Metro monopoliza los diez primeros puestos pese a concentrar
menos del 2\,\% de las aristas de servicio real: su ventaja no es
volumétrica sino posicional ---sus estaciones ocupan los puntos donde
más rutas se cruzan.

La distribución geográfica del top~10 es igualmente estrecha: siete
de los diez nodos quedan en Benito Juárez o Cuauhtémoc, lo que revela
que la intermediación de la red está concentrada en un corredor
reducido del centro-poniente de la cuenca.

La Figura~\ref{fig:grafica-bv-nodos} extiende ese ranking a los
principales nodos de la red.

\begin{figure}[H]
    \centering
    \includegraphics[width=0.92\textwidth]{%
        Figures/Mapas/Tableau_Resources/RedActual/VFTModel/UniGrafica/%
        Grafica_VFTModel_2026_Centralidad_Intermediacion_Nodos}
    \caption[Centralidad de intermediación $B(v)$ --- top nodos de la red]{%
    Ranking de los nodos con mayor centralidad de intermediación
    ($B(v)$) en la red de transporte de la \zmvm{}.
    Corrida: 2026-09-02.}
    \label{fig:grafica-bv-nodos}
    \fuentefigura{Elaboración propia con \textbf{Tableau};
    datos: \texttt{VFTModel} \citep{cabrera2026VFTModel}.}
\end{figure}

La centralidad topológica no siempre coincide con la afluencia real.
Para verificarlo, la Figura~\ref{fig:grafica-top10-estaciones} muestra
las diez estaciones con mayor volumen de usuarios según los registros
de Apímetro ---un contraste directo con el ranking de $B(v)$ del
Cuadro~\ref{tab:bv-top10}.

\begin{figure}[H]
    \centering
    \includegraphics[width=0.92\textwidth]{%
        Figures/Mapas/Tableau_Resources/RedActual/Apimetro/UniGrafica/%
        Grafica_Apimetro_2026_Top_10_estaciones}
    \caption[Top 10 estaciones por afluencia de usuarios --- \zmvm{}]{%
    Diez estaciones con mayor volumen de usuarios registrado por
    \texttt{Apímetro} en la red de transporte de la \zmvm{}.
    Corrida: 2026-04-18.}
    \label{fig:grafica-top10-estaciones}
    \fuentefigura{Elaboración propia con \textbf{Tableau};
    datos: \texttt{Apímetro} \citep{cabrera2025apimetro}.}
\end{figure}

La Figura~\ref{fig:grafica-bv-sistema} agrega esos valores por sistema,
mostrando cuánto peso acumulado concentra cada modo en la intermediación
global de la red.

\begin{figure}[H]
    \centering
    \includegraphics[width=0.92\textwidth]{%
        Figures/Mapas/Tableau_Resources/RedActual/VFTModel/UniGrafica/%
        Grafica_VFTModel_2026_Centralidad_Intermediacion_sistema}
    \caption[Centralidad de intermediación $B(v)$ acumulada por sistema]{%
    Suma acumulada de $B(v)$ agrupada por sistema de transporte.
    El Metro concentra la mayor proporción de intermediación global
    de la red. Corrida: 2026-09-02.}
    \label{fig:grafica-bv-sistema}
    \fuentefigura{Elaboración propia con \textbf{Tableau};
    datos: \texttt{VFTModel} \citep{cabrera2026VFTModel}.}
\end{figure}

Para traducirlo al lenguaje urbano: imaginemos que cada par de colonias
de la \zmvm{} traza su ruta más rápida hacia el destino. Las estaciones
que aparecen con mayor frecuencia en ese ejercicio colectivo no son
necesariamente las más concurridas en términos de afluencia, sino las
más difíciles de sustituir: su ausencia obligaría a miles de trayectos
a buscar desvíos más largos o, en el peor caso, a no encontrarlos.
$B(v)$ mide exactamente esa indispensabilidad estructural.

La Figura~\ref{fig:mapa-betweenness} traslada ese ranking al territorio,
mostrando cómo la intermediación topológica se distribuye espacialmente
sobre la cuenca metropolitana.

\begin{figure}[H]
    \centering
    \includegraphics[width=0.92\textwidth]{%
        Figures/Mapas/Qgis_Resources/RedActual/Mapa_Betweenness}
    \caption[Mapa de centralidad de intermediación $B(v)$ --- \zmvm{}]{%
    Distribución espacial de la centralidad de intermediación ($B(v)$)
    en la red de transporte de la \zmvm{}. El corredor de alta
    intermediación traza una banda en el centro-poniente de la cuenca.
    Corrida: 2026-09-02.}
    \label{fig:mapa-betweenness}
    \fuentefigura{Elaboración propia con \textbf{QGIS};
    datos: \texttt{VFTModel} \citep{cabrera2026VFTModel}.}
\end{figure}

La concentración que revelan los datos tabulares adquiere forma geográfica
en el mapa: el corredor de alta intermediación no cubre la metrópoli de
manera uniforme, sino que traza una banda relativamente estrecha sobre el
centro-poniente, dejando amplias zonas periféricas con valores bajos o
nulos. Esas zonas ---las mismas que registraron menor cobertura y mayor
Factor de Desviación en Fase~1--- son también las que dependen de los
nodos del corredor central para cualquier viaje de largo aliento.

La Figura~\ref{fig:dashboard-bv} consolida la visión integral del indicador
sobre el conjunto de la red.

\begin{figure}[H]
    \centering
    \includegraphics[width=0.92\textwidth]{%
        Figures/Mapas/Tableau_Resources/RedActual/VFTModel/Dashboards/%
        Dashboard_VFTModel_2026_Centralidad_intermediacion}
    \caption[Dashboard integral de centralidad de intermediación $B(v)$ --- \zmvm{}]{%
    Vista integral del indicador de centralidad de intermediación ($B(v)$)
    sobre la red de transporte de la \zmvm{}.
    Corrida: 2026-09-02.}
    \label{fig:dashboard-bv}
    \fuentefigura{Elaboración propia con \textbf{Tableau};
    datos: \texttt{VFTModel} \citep{cabrera2026VFTModel}.}
\end{figure}

Los tres indicadores de Fase~3 ---tiempo de viaje, intermediación y
vulnerabilidad preliminar--- apuntan en la misma dirección: la red
concentra su capacidad de articulación en un corredor reducido del
centro-poniente, mientras la periferia sur y oriente queda
estructuralmente subordinada a él. La sección siguiente examina qué
ocurre cuando alguno de esos nodos críticos desaparece de la red.

% !TeX root = ../../tesis.tex
%%%%%%%%%%%%%%%%%%%%%%%%%%%%%%%%%%%%%%%%%%%%%%%%%%%%%%%%%%%%%%%%%%%
% SECCIÓN 5.5.3 — VULNERABILIDAD ESTRUCTURAL ΔE
%%%%%%%%%%%%%%%%%%%%%%%%%%%%%%%%%%%%%%%%%%%%%%%%%%%%%%%%%%%%%%%%%%%
\subsection{Vulnerabilidad Estructural (\texorpdfstring{$\Delta E$}{Delta E})}
\label{subsec:vulnerabilidad-estructural}

Los resultados de esta subsección tienen carácter preliminar. La
simulación de vulnerabilidad estructural se realizó sobre los nodos
del grafo construido por \texttt{VFTModel}, pero no fue ejecutada
sobre la capa de red integrada de \texttt{Transport-gis-mjg}; en
consecuencia, los valores reportados no han sido verificados contra
la representación completa del sistema. Se presentan como primer
acercamiento al indicador $\Delta E$, con la intención de enmarcar
el método y señalar las líneas de trabajo futuro.

A partir del valor de referencia $T = 108.92$~min
(§\ref{subsec:tiempo-viaje-promedio}), el experimento calcula la
caída relativa de eficiencia $\Delta E$ cuando se retira cada uno de
los cinco nodos de mayor $B(v)$ (§\ref{subsec:robustez-geometrica}).
La eficiencia base es $E_0 = 1/108.92$; tras remover el nodo $v$,
se recalcula el tiempo promedio $T_f$ y se obtiene
$\Delta E = (E_0 - E_f)/E_0$.

El Cuadro~\ref{tab:delta-e-top5} registra los resultados del
experimento sobre los cinco nodos seleccionados.

\begin{table}[H]
    \centering
    \caption[Impacto de remoción de nodos críticos --- $\Delta E$ exploratorio]{%
    Caída relativa de eficiencia $\Delta E$ al remover cada nodo del
    componente gigante. $E_0 = 1/108.92$; $E_f = 1/T_f$. Corrida
    exploratoria: 2026-09-02.}
    \label{tab:delta-e-top5}
    \small
    \begin{tabular}{|r|l|r|r|r|}
        \hline
        \textbf{\#} & \textbf{Nodo removido} &
        \textbf{$B(v)$} & \textbf{$T_f$ (min)} & \textbf{$\Delta E$} \\
        \hline
        1 & Tacubaya        & 0.2179 & 111.80 & 2.58\,\% \\
        2 & Mixcoac         & 0.1725 & 110.19 & 1.15\,\% \\
        3 & Hidalgo         & 0.1595 & 109.90 & 0.89\,\% \\
        4 & Balderas        & 0.1433 & 109.57 & 0.59\,\% \\
        5 & Lázaro Cárdenas & 0.1376 & 110.13 & 1.10\,\% \\
        \hline
    \end{tabular}
    \fuentefigura{Elaboración propia con \texttt{VFTModel}
    \citep{cabrera2026VFTModel}. Corrida exploratoria: 2026-09-02.}
\end{table}

Tacubaya encabeza tanto el ranking de $B(v)$ como el de $\Delta E$: su
remoción eleva el tiempo promedio de la red en 2.88~min, traducido en
una caída de eficiencia del 2.58\,\%. El dato es coherente con la
posición estructural del nodo ---cruce de las Líneas 1, 7 y 9 del
Metro--- y su papel de puente entre el corredor poniente y el centro.
Los cuatro nodos restantes producen caídas entre 0.59\,\% y 1.15\,\%,
lo que confirma que ningún nodo individual colapsa la red, pero la
pérdida acumulada de varios podría alcanzar umbrales críticos.

Los tres indicadores de Fase~3 convergen en un mismo diagnóstico
estructural. $T = 108.92$~min revela que el tiempo de viaje promedio
entre cualquier par de la \zmvm{} supera con amplitud los umbrales de
movilidad cotidiana. $B(v)$ identifica el corredor que hace posible
ese tiempo: un conjunto reducido de estaciones del Metro en el eje
centro-poniente concentra la mayor parte de los caminos mínimos de la
red. Y la simulación exploratoria de $\Delta E$ confirma que la
remoción de esos nodos, aunque no colapsa la red, la degrada de
forma mensurable ---con Tacubaya produciendo el mayor impacto
individual registrado.

El patrón es consistente con el hallazgo de Fase~1: las zonas que
registraron menor cobertura y mayor Factor de Desviación son también
las que más dependen del corredor central para cualquier trayecto de
largo aliento. La siguiente sección integra estos indicadores en un
diagnóstico territorial que localiza las zonas de déficit estructural
de la red.


% \input{5-Analisis-Red-Actual/5-6-Diagnostico-Territorial}
